\begin{tikzpicture}[font=\small,>=stealth]
\draw[rounded corners=3pt,gray!45] (0,0) rectangle (14.2,12);

\node at (7.1,11.45) {先看熟悉的投影：斜边乘余弦，就是水平位移};
\begin{scope}[shift={(5.2,6.1)},scale=1.55]
\draw[gray] (-2.1,0)--(1.65,0);
\draw[blue!75!black,very thick] (0,0)--(.6428571429,1.1134612334) node[midway,left=8pt] {$r$};
\draw[orange!85!black,very thick] (0,0)--(-1.5,-2.5980762114) node[midway,right=8pt] {$s$};
\draw[gray,dashed] (.6428571429,1.1134612334)--(.6428571429,0);
\draw[gray,dashed] (-1.5,-2.5980762114)--(-1.5,0);
\draw[blue!65!black,thick] (0,0)--(.6428571429,0);
\draw[orange!85!black,thick] (-1.5,0)--(0,0);
\draw[gray] (.5,0)--(.5,.14)--(.6428571429,.14);
\draw[gray] (-1.35,0)--(-1.35,-.15)--(-1.5,-.15);
\draw[violet] (.45,0) arc[start angle=0,end angle=60,radius=.45];
\node[violet] at (.83,.3) {$\theta$};
\fill (0,0) circle(2pt) node[below right=6pt] {$P$};
\fill (.6428571429,1.1134612334) circle(2pt) node[above=8pt] {$A$};
\fill (-1.5,-2.5980762114) circle(2pt) node[left=10pt] {$B$};
\node[above=9pt,orange!85!black] at (-.75,0) {$s\cos\theta$};
\node[below right=8pt,blue!75!black] at (.6428571429,0) {$r\cos\theta$};
\end{scope}
\node[anchor=west,align=left] at (9.1,7.2)
 {图中为锐角\\[14pt]从 $P$ 到 $A$：向右\\[12pt]水平位移 $+r\cos\theta$\\[18pt]
 从 $P$ 到 $B$：向左\\[12pt]水平位移 $-s\cos\theta$};
\node at (7.1,1.2) {先只认“走多远”和“横向走多远”，不需要椭圆的新知识。};
\node at (7.1,.5) {钝角时余弦为负，位移式仍成立；不要把位移的负号当成长度为负。};
\end{tikzpicture}
